\begin{pastexam}[label={px:ntu-107-ie-ds-13}]{臺大工工所 107 學年度 第 13 題}
（是非；命題正確填 (A)，錯誤填 (B)；答對得 4 分，答錯倒扣 4 分，至多扣到該大題零分）

In a red-black tree, every subtree is also a red-black tree.
\end{pastexam}

\begin{pastsoln}
答案是 (B)，錯誤。

紅黑樹的性質（資結書第~10.3~節）有一條是\textbf{根是黑色}。 紅黑樹裡常有紅色的節點，以紅色節點
為根的子樹，根是紅的，違反這一條。

反例：鍵 $1, 2, 3$，根 $2$ 是黑色，$1$、$3$ 是紅色。 這是合法的紅黑樹（每條到 \NIL 的路徑都經過
一個內部黑色節點）。 以 $1$ 為根的子樹只有一個紅色節點，根不是黑的，不是紅黑樹。

其他幾條性質（紅色節點的子節點是黑的、每個節點到底下各 \NIL 的黑高度相同）只看節點底下的東西，
在子樹裡仍然成立；所以子樹只差在根的顏色，把根改塗黑就是紅黑樹。 但題目問的是子樹本身，答案是錯。
對照第 10 題，AVL 樹的條件全部只看節點底下，子樹一定還是 AVL 樹。
\end{pastsoln}
